\frfilename{mt243.tex}
\versiondate{30.4.04}
     
\def\chaptername{Ruang fungsi}
\def\sectionname{$L^{\infty}$}
     
\newsection{243}
     
Ruang Banach klasik kedua dalam teori ukur yang saya bahas adalah ruang
$L^{\infty}$. Sebagaimana akan tampak di bawah, $L^{\infty}$ merupakan
pasangan polar dari $L^1$, yakni lawan yang terkait; untuk ruang ukur
`biasa', ruang ini sebenarnya merupakan dual dari $L^1$ (243F-243G).
     
\leader{243A}{Definisi} Misalkan $(X,\Sigma,\mu)$ sebarang ruang ukur.
Misalkan $\eusm L^{\infty}=\eusm L^{\infty}(\mu)$ himpunan fungsi
$f\in\eusm L^0=\eusm L^0(\mu)$ yang {\bf terbatas secara esensial}, yaitu
sedemikian sehingga terdapat $M\ge 0$ dengan
$\{x:x\in\dom f,\,|f(x)|\le M\}$ koterabaikan, dan tuliskan
     
\Centerline{$L^{\infty}=L^{\infty}(\mu)
=\{f^{\ssbullet}:f\in\eusm L^{\infty}(\mu)\}\subseteq L^0(\mu)$.}
     
Perhatikan bahwa\cmmnt{ jika $f\in\eusm L^{\infty}$,
$g\in\eusm L^0$ dan $g\eae f$, maka $g\in\eusm L^{\infty}$; dengan
demikian}
$\eusm L^{\infty}=\{f:f\in\eusm L^0,\,f^{\ssbullet}\in L^{\infty}\}$.
     
\vleader{72pt}{243B}{Teorema} Misalkan $(X,\Sigma,\mu)$ sebarang ruang ukur. Maka
     
(a) $L^{\infty}=L^{\infty}(\mu)$ merupakan subruang linear dari
$L^0=L^0(\mu)$.
     
(b) Jika $u\in L^{\infty}$, $v\in L^0$ dan $|v|\le|u|$, maka
$v\in L^{\infty}$. Akibatnya, $|u|$, $u\vee v$, $u\wedge v$,
$u^+\cmmnt{\mskip5mu =u\vee 0}$ dan $u^-\cmmnt{\mskip5mu =(-u)\vee 0}$
termasuk dalam $L^{\infty}$ untuk semua $u$, $v\in L^{\infty}$.
     
(c) Dengan menuliskan $e=\chi X^{\ssbullet}$, yaitu kelas ekuivalensi
dalam $L^0$ dari fungsi konstan yang bernilai $1$, suatu unsur $u$ dari
$L^0$ termasuk dalam $L^{\infty}$ jika dan hanya jika terdapat $M\ge 0$
sedemikian sehingga $|u|\le Me$.
     
(d) Jika $u$, $v\in L^{\infty}$, maka $u\times v\in L^{\infty}$.
     
(e) Jika $u\in L^{\infty}$ dan $v\in L^1=L^1(\mu)$, maka
$u\times v\in L^1$.
     
\proof{{\bf (a)} Jika $f$, $g\in\eusm
L^{\infty}=\eusm L^{\infty}(\mu)$ dan $c\in\Bbb R$, maka $f+g$,
$cf\in\eusm L^{\infty}$.
\Prf\ Kita mempunyai $M_1$, $M_2\ge 0$ sedemikian sehingga
$|f|\le M_1$ hampir di mana-mana dan $|g|\le M_2$ hampir di mana-mana.
Sekarang
     
\Centerline{$|f+g|\le|f|+|g|\le M_1+M_2$ hampir di mana-mana,
\quad $|cf|\le|c|M_1$ hampir di mana-mana,}
     
\noindent sehingga $f+g$, $cf\in\eusm L^{\infty}$.
\ \QeD\ Segera diperoleh bahwa $u+v$, $cu\in L^{\infty}$ setiap kali
$u$, $v\in L^{\infty}$ dan $c\in\Bbb R$.
     
\medskip
     
{\bf (b)(i)} Ambil $f\in\eusm L^{\infty}$,
$g\in\eusm L^0=\eusm L^0(\mu)$ sedemikian sehingga
$u=f^{\ssbullet}$ dan $v=g^{\ssbullet}$. Maka $|g|\leae|f|$.
Misalkan $M\ge 0$ sedemikian sehingga $|f|\le M$ hampir di mana-mana;
maka $|g|\le M$ hampir di mana-mana, sehingga
$g\in\eusm L^{\infty}$ dan $v\in L^{\infty}$.
     
\medskip
     
\quad{\bf (ii)} Sekarang $|\,|u|\,|=|u|$, sehingga
$|u|\in L^{\infty}$ setiap kali $u\in L^{\infty}$. Selain itu,
$u\vee v=\bover12(u+v+|u-v|)$ dan
$u\wedge v=\bover12(u+v-|u-v|)$ termasuk dalam $L^{\infty}$ untuk semua
$u$, $v\in L^{\infty}$.
     
\medskip
     
{\bf (c)}(i) Jika $u\in L^{\infty}$, ambil
$f\in\eusm L^{\infty}$ sedemikian sehingga $f^{\ssbullet}=u$.
Maka terdapat $M\ge 0$ sedemikian sehingga $|f|\le M$ hampir di
mana-mana, sehingga $|f|\leae M\chi X$ dan $|u|\le Me$. (ii) Tentu saja
$\chi X\in\eusm L^{\infty}$, sehingga $e\in L^{\infty}$, dan jika
$u\in L^0$ serta $|u|\le Me$, maka $u\in L^{\infty}$ menurut (b).
     
\medskip
     
{\bf (d)} $f\times g\in\eusm L^{\infty}$ setiap kali $f$,
$g\in\eusm L^{\infty}$. \Prf\ Jika $|f|\le M_1$ hampir di mana-mana
dan $|g|\le M_2$ hampir di mana-mana, maka
     
\Centerline{$|f\times g|=|f|\times|g|\le M_1M_2$ hampir di mana-mana.
\Qed}
     
\noindent Jadi $u\times v\in L^{\infty}$ untuk semua $u$,
$v\in L^{\infty}$.
     
\medskip
     
{\bf (e)} Jika $f\in\eusm L^{\infty}$ dan
$g\in\eusm L^1=\eusm L^1(\mu)$, maka terdapat $M\ge 0$ sedemikian
sehingga $|f|\le M$ hampir di mana-mana, sehingga
$|f\times g|\leae M|g|$; karena $M|g|$ terintegralkan dan $f\times g$
terukur secara virtual, $f\times g$ terintegralkan dan $u\times v\in L^1$.
}%end of proof of 243B
     
\leader{243C}{Struktur terurut $L^{\infty}$}
Misalkan $(X,\Sigma,\mu)$ sebarang ruang ukur. Maka
$L^{\infty}=L^{\infty}(\mu)$, sebagai subruang linear dari
$L^0=L^0(\mu)$, mewarisi suatu urutan parsial yang menjadikannya ruang
linear terurut parsial\cmmnt{ (bandingkan 242Ca)}. Karena
$|u|\in L^{\infty}$ setiap kali $u\in L^{\infty}$\cmmnt{ (243Bb)},
$u\wedge v$ dan $u\vee v$ termasuk dalam $L^{\infty}$ setiap kali $u$,
$v\in L^{\infty}$, dan $L^{\infty}$ merupakan ruang Riesz\cmmnt{
(bandingkan 242Cd)}.
     
\cmmnt{Perilaku }$L^{\infty}$\cmmnt{ sebagai ruang Riesz didominasi
oleh fakta bahwa ruang ini} mempunyai suatu {\bf satuan urutan} $e$
dengan sifat bahwa
     
\Centerline{untuk setiap $u\in L^{\infty}$ terdapat $M\ge 0$ sedemikian
sehingga $|u|\le Me$\dvro{.}{}}

\cmmnt{\noindent (243Bc).}
     
\leader{243D}{Norma $L^{\infty}$} Misalkan $(X,\Sigma,\mu)$ sebarang
ruang ukur.
     
\header{243Da}{\bf (a)} Untuk
$f\in\eusm L^{\infty}=\eusm L^{\infty}(\mu)$,\cmmnt{ katakanlah}
{\bf supremum esensial} dari $|f|$ adalah
     
\Centerline{$\esssup|f|=
\inf\{M:M\ge 0,\,\{x:x\in\dom f,\,|f(x)|\le M\}$ koterabaikan$\}$.}
     
\noindent Maka $|f|\le\esssup|f|$ hampir di mana-mana.
\prooflet{\Prf\ Tetapkan $M=\esssup|f|$. Untuk setiap $n\in\Bbb N$,
terdapat $M_n\le M+2^{-n}$ sedemikian sehingga $|f|\le M_n$ hampir di
mana-mana. Sekarang
     
\Centerline{$\{x:|f(x)|\le M\}=\bigcap_{n\in\Bbb N}\{x:|f(x)|\le M_n\}$}
     
\noindent koterabaikan, sehingga $|f|\le M$ hampir di mana-mana.\ \Qed}
     
\header{243Db}{\bf (b)} Jika $f$, $g\in\eusm L^{\infty}$ dan
$f\eae g$, maka $\esssup|f|=\esssup|g|$. Dengan demikian kita dapat
mendefinisikan suatu fungsional $\|\,\|_{\infty}$ pada
$L^{\infty}=L^{\infty}(\mu)$ dengan menetapkan
$\|u\|_{\infty}=\esssup|f|$ setiap kali $u=f^{\ssbullet}$.
     
\medskip
     
\spheader 243Dc\dvro{ Untuk}{ Dari (a), kita melihat bahwa, untuk}
sebarang $u\in L^{\infty}$,
$\|u\|_{\infty}=\min\{\gamma:|u|\le\gamma e\}$, dengan\cmmnt{, seperti
sebelumnya,} $e=\chi X^{\ssbullet}\in L^{\infty}$. \cmmnt{Akibatnya}
$\|\,\|_{\infty}$ merupakan suatu norma pada $L^{\infty}$.
\prooflet{\Prf(i) Jika $u$, $v\in L^{\infty}$, maka
     
     
\Centerline{$|u+v|\le|u|+|v|\le(\|u\|_{\infty}+\|v\|_{\infty})e$}
     
\noindent sehingga
$\|u+v\|_{\infty}\le\|u\|_{\infty}+\|v\|_{\infty}$. (ii) Jika
$u\in L^{\infty}$ dan $c\in\Bbb R$, maka
     
\Centerline{$|cu|=|c||u|\le|c|\|u\|_{\infty}e$,}
     
\noindent sehingga $\|cu\|_{\infty}\le|c|\|u\|_{\infty}$. (iii) Jika
$\|u\|_{\infty}=0$, terdapat $f\in\eusm L^{\infty}$ sedemikian sehingga
$f^{\ssbullet}=u$ dan $|f|\le\|u\|_{\infty}$ hampir di mana-mana;
sekarang $f=0$ hampir di mana-mana, sehingga $u=0$. \Qed}
     
\spheader 243Dd Perhatikan pula bahwa jika $u\in L^0$,
$v\in L^{\infty}$ dan $|u|\le|v|$, maka
$|u|\le\|v\|_{\infty}e$, sehingga $u\in L^{\infty}$ dan
$\|u\|_{\infty}\le\|v\|_{\infty}$; demikian pula,
     
\Centerline{$\|u\times v\|_{\infty}\le\|u\|_{\infty}\|v\|_{\infty}$,
\quad$\|u\vee v\|_{\infty}\le\max(\|u\|_{\infty},\|v\|_{\infty})$}
     
\noindent untuk semua $u$, $v\in L^{\infty}$. \cmmnt{Dengan demikian}
$L^{\infty}$ merupakan aljabar Banach komutatif\cmmnt{ (2A4J)}.
     
\spheader 243De \cmmnt{Selain itu,}

\Centerline{$|\int u\times v|\le\int|u\times v|
=\|u\times v\|_1\le\|u\|_1\|v\|_{\infty}$}
     
\noindent setiap kali $u\in L^1$ dan $v\in L^{\infty}$\dvro{.}{,
karena
     
\Centerline{$|u\times v|=|u|\times|v|\le|u|\times
\|v\|_{\infty}e=\|v\|_{\infty}|u|$.}
}
     
\spheader 243Df
Perhatikan bahwa jika $u$, $v$ merupakan anggota nonnegatif dari
$L^{\infty}$, maka
     
\Centerline{$\|u\vee v\|_{\infty}
=\max(\|u\|_{\infty},\|v\|_{\infty})$\dvro{.}{;}}
     
\cmmnt{\noindent ini karena, untuk sebarang $\gamma\ge 0$,
     
\Centerline{$u\vee v\le\gamma e\,\iff\,u\le\gamma e$ dan
$v\le\gamma e$.}
}
     
\leader{243E}{Teorema} Untuk sebarang ruang ukur $(X,\Sigma,\mu)$,
$L^{\infty}=L^{\infty}(\mu)$ merupakan kisi Banach di bawah
$\|\,\|_{\infty}$.
     
\proof{{\bf (a)} Kita telah mengetahui bahwa
$\|u\|_{\infty}\le\|v\|_{\infty}$ setiap kali $|u|\le|v|$ (243Dd);
jadi kita hanya perlu memeriksa bahwa $L^{\infty}$ lengkap di bawah
$\|\,\|_{\infty}$. Misalkan $\sequencen{u_n}$ suatu barisan Cauchy dalam
$L^{\infty}$. Untuk setiap $n\in\Bbb N$, pilih
$f_n\in\eusm L^{\infty}=\eusm L^{\infty}(\mu)$ sedemikian sehingga
$f_n^{\ssbullet}=u_n$ dalam $L^{\infty}$. Untuk semua $m$,
$n\in\Bbb N$, $(f_m-f_n)^{\ssbullet}=u_m-u_n$. Akibatnya
     
\Centerline{$E_{mn}=\{x:|f_m(x)-f_n(x)|>\|u_m-u_n\|_{\infty}\}$}
     
\noindent terabaikan, menurut 243Da. Ini berarti bahwa
     
\Centerline{$E
=\bigcap_{n\in\Bbb N}\{x:x\in\dom f_n,\,|f_n(x)|\le\|u_n\|_{\infty}\}
  \setminus\bigcup_{m,n\in\Bbb N}E_{mn}$}
     
\noindent koterabaikan. Namun, untuk setiap $x\in E$,
$|f_m(x)-f_n(x)|\le\|u_m-u_n\|_{\infty}$ untuk semua $m$,
$n\in\Bbb N$, sehingga $\sequencen{f_n(x)}$ merupakan barisan Cauchy
dengan limit dalam $\Bbb R$. Jadi $f=\lim_{n\to\infty}f_n$ terdefinisi
hampir di mana-mana. Selain itu, setidaknya untuk $x\in E$,
     
\Centerline{$|f(x)|\le\sup_{n\in\Bbb N}\|u_n\|_{\infty}<\infty$,}
     
\noindent sehingga $f\in\eusm L^{\infty}$ dan
$u=f^{\ssbullet}\in L^{\infty}$. Jika $m\in\Bbb N$, maka, untuk setiap
$x\in E$,

\Centerline{$|f(x)-f_m(x)|\le\sup_{n\ge m}|f_n(x)-f_m(x)|
\le\sup_{n\ge m}\|u_n-u_m\|_{\infty}$,}
     
\noindent sehingga
     
\Centerline{$\|u-u_m\|_{\infty}\le\sup_{n\ge m}\|u_n-u_m\|_{\infty}
\to 0$}
     
\noindent ketika $m\to\infty$, dan $u=\lim_{m\to\infty}u_m$ dalam
$L^{\infty}$. Karena $\sequencen{u_n}$ sembarang, $L^{\infty}$ lengkap.
}
     
\leader{243F}{Dualitas antara $L^{\infty}$ dan $L^1$}
Misalkan $(X,\Sigma,\mu)$ sebarang ruang ukur.
     
     
\cmmnt{\header{243Fa}{\bf (a)}
Saya telah menyatakan bahwa jika $u\in L^{1}=L^{1}(\mu)$ dan
$v\in L^{\infty}=L^{\infty}(\mu)$, maka $u\times v\in L^1$ dan
$|\int u\times v|\le\|u\|_1\|v\|_{\infty}$ (243Bd, 243De).
}
     
\header{243Fb}{\bf (b)} \dvro{Kita}{Akibatnya kita} mempunyai suatu
operator linear terbatas $T$ dari $L^{\infty}$ ke dual ruang bernorma
$(L^1)^*$ dari $L^1$, yang diberikan dengan menuliskan
     
\Centerline{$(Tv)(u)=\int u\times v$ untuk semua $u\in L^1$,
$v\in L^{\infty}$.}

\prooflet{\noindent\Prf\ (i)
Menurut (a), $(Tv)(u)$ terdefinisi dengan baik untuk $u\in L^1$,
$v\in L^{\infty}$. (ii) Jika $v\in L^{\infty}$, $u$, $u_1$,
$u_2\in L^1$ dan $c\in\Bbb R$, maka
     
$$\eqalign{(Tv)(u_1+u_2)
&=\int(u_1+u_2)\times v
=\int(u_1\times v)+(u_2\times v)\cr
&=\int u_1\times v+\int u_2\times v
=(Tv)(u_1)+(Tv)(u_2),\cr}$$
     
\Centerline{$(Tv)(cu)=\int cu\times v=\int c(u\times v)
=c\int u\times v=c(Tv)(u)$.}
     
\noindent Ini menunjukkan bahwa $Tv:L^1\to\Bbb R$ merupakan fungsional
linear untuk setiap $v\in L^{\infty}$. (iii) Selanjutnya, untuk sebarang
$u\in L^1$ dan $v\in L^{\infty}$,
     
\Centerline{$|(Tv)(u)|=|\int u\times v|\le\|u\times v\|_1
\le\|u\|_1\|v\|_{\infty}$,}
     
\noindent sebagaimana dinyatakan dalam (a). Ini berarti bahwa
$Tv\in (L^1)^*$ dan $\|Tv\|\le\|v\|_{\infty}$ untuk setiap
$v\in L^{\infty}$. (iv) Jika $v$, $v_1$, $v_2\in L^{\infty}$,
$u\in L^1$ dan $c\in\Bbb R$, maka
     
$$\eqalign{T(v_1+v_2)(u)
&=\int(v_1+v_2)\times u
=\int(v_1\times u)+(v_2\times u)\cr
&=\int v_1\times u+\int v_2\times u
=(Tv_1)(u)+(Tv_2)(u)\cr
&=(Tv_1+Tv_2)(u),\cr}$$
     
\Centerline{$T(cv)(u)=\int cv\times u=c\int v\times u
=c(Tv)(u)=(cTv)(u)$.}
     
\noindent Karena $u$ sembarang,
$T(v_1+v_2)=Tv_1+Tv_2$ dan $T(cv)=c(Tv)$; jadi
$T:L^{\infty}\to(L^1)^*$ linear. (v) Dengan mengingat dari (iii) bahwa
$\|Tv\|\le\|v\|_{\infty}$ untuk setiap $v\in L^{\infty}$, kita melihat
bahwa $\|T\|\le 1$.\ \Qed}
     
\header{243Fc}{\bf (c)} \dvro{Kita}{Argumen yang persis sama menunjukkan
bahwa kita} mempunyai suatu operator linear $T':L^1\to (L^{\infty})^*$,
yang diberikan dengan menuliskan
     
\Centerline{$(T'u)(v)=\int u\times v$ untuk semua $u\in L^1$,
$v\in L^{\infty}$,}
     
\noindent dan\cmmnt{ bahwa} $\|T'\|$ juga paling besar $1$.

\leader{243G}{Teorema} Misalkan $(X,\Sigma,\mu)$ suatu ruang ukur, dan
$T:L^{\infty}(\mu)\to(L^1(\mu))^*$ peta kanonik yang diuraikan dalam
243F. Maka
     
(a) $T$ injektif jika dan hanya jika $(X,\Sigma,\mu)$ semihingga, dan
dalam hal ini peta tersebut mempertahankan norma;
     
(b) $T$ bijektif jika dan hanya jika $(X,\Sigma,\mu)$ dapat dilokalkan,
dan dalam hal ini peta tersebut merupakan isomorfisme ruang bernorma.
     
\proof{{\bf (a)(i)} Andaikan $T$ injektif, dan $E\in\Sigma$ mempunyai
$\mu E=\infty$. Maka $\chi E$ tidak sama hampir di mana-mana dengan
$\tbf{0}$, sehingga $(\chi E)^{\ssbullet}\ne 0$ dalam $L^{\infty}$, dan
$T(\chi E)^{\ssbullet}\ne 0$; misalkan $u\in L^1$ sedemikian sehingga
$T(\chi E)^{\ssbullet}(u)\ne 0$, yaitu
$\int u\times(\chi E)^{\ssbullet}\ne 0$. Nyatakan $u$ sebagai
$f^{\ssbullet}$ dengan $f$ terintegralkan; maka $\int_Ef\ne 0$, sehingga
$\int_E|f|\ne 0$. Misalkan $g$ suatu fungsi sederhana sedemikian sehingga
$0\le g\leae|f|$ dan $\int g>\int|f|-\int_E|f|$; maka
$\int_Eg\ne 0$. Nyatakan $g$ sebagai $\sum_{i=0}^na_i\chi E_i$ dengan
$\mu E_i<\infty$ untuk setiap $i$; maka
$0\ne\sum_{i=0}^na_i\mu(E_i\cap E)$, sehingga terdapat $i\le n$
sedemikian sehingga $\mu(E\cap E_i)\ne 0$, dan sekarang $E\cap E_i$
merupakan himpunan bagian terukur dari $E$ yang berukuran hingga dan
tak nol.
     
Karena $E$ sembarang, ini menunjukkan bahwa $(X,\Sigma,\mu)$ harus
semihingga jika $T$ injektif.
     
\medskip
     
\quad{\bf (ii)} Sekarang andaikan $(X,\Sigma,\mu)$ semihingga, dan
$v\in L^{\infty}$ tak nol. Nyatakan $v$ sebagai $g^{\ssbullet}$ dengan
$g:X\to\Bbb R$ terukur; maka $g\in\eusm L^{\infty}$. Ambil sebarang
$a\in\ooint{0,\|v\|_{\infty}}$; maka
$E=\{x:|g(x)|\ge a\}$ mempunyai ukuran tak nol. Misalkan $F\subseteq E$
suatu himpunan terukur yang berukuran hingga dan tak nol, dan tetapkan
$f(x)=|g(x)|/g(x)$ jika $x\in F$, serta $0$ selainnya; maka
$f\in\eusm L^1$ dan $(f\times g)(x)\ge a$ untuk $x\in F$, sehingga,
dengan menetapkan $u=f^{\ssbullet}\in L^1$, kita mempunyai
     
\Centerline{$(Tv)(u)=\int u\times v=\int f\times g\ge a\mu F=a\int|f|
=a\|u\|_1>0$.}
     
\noindent Ini menunjukkan bahwa $\|Tv\|\ge a$; karena $a$ sembarang,
$\|Tv\|\ge\|v\|_{\infty}$. Kita telah mengetahui dari 243F bahwa
$\|Tv\|\le\|v\|_{\infty}$, sehingga
$\|Tv\|=\|v\|_{\infty}$ untuk setiap $v\in L^{\infty}$ yang tak nol;
hal yang sama tentu benar untuk $v=0$, sehingga $T$ mempertahankan norma
dan injektif.
     
\medskip
     
{\bf (b)(i)} Dengan menggunakan (a) dan definisi `dapat dilokalkan',
kita melihat bahwa di bawah salah satu syarat yang diajukan,
$(X,\Sigma,\mu)$ semihingga dan $T$ injektif serta mempertahankan norma.
Karena itu saya hanya perlu menunjukkan bahwa peta tersebut surjektif jika dan
hanya jika $(X,\Sigma,\mu)$ dapat dilokalkan.
     
\medskip
     
\quad{\bf (ii)} Andaikan $T$ surjektif dan $\Cal E\subseteq\Sigma$.
Misalkan $\Cal F$ keluarga gabungan hingga dari anggota-anggota
$\Cal E$, dengan $\emptyset$ dihitung sebagai gabungan tanpa anggota
$\Cal E$, sehingga $\Cal F$ tertutup terhadap gabungan hingga dan,
untuk setiap $G\in\Sigma$, $E\setminus G$ terabaikan untuk setiap
$E\in\Cal E$ jika dan hanya jika $E\setminus G$ terabaikan untuk setiap
$E\in\Cal F$.
     
Jika $u\in L^1$, maka
$h(u)=\lim_{E\in\Cal F,E\uparrow}\int_Eu$ ada dalam $\Bbb R$.
\prooflet{\Prf\ Jika $u$ nonnegatif, maka
     
     
\Centerline{$h(u)=\sup\{\int_Eu:E\in\Cal F\}\le\int u<\infty$.}
     
\noindent Untuk $u$ yang lain, kita dapat menyatakan $u$ sebagai
$u_1-u_2$, dengan $u_1$ dan $u_2$ nonnegatif, dan sekarang
$h(u)=h(u_1)-h(u_2)$.\ \Qed}
     
Jelas bahwa $h:L^1\to\Bbb R$ linear, karena merupakan limit dari
fungsional-fungsional linear $u\mapsto\int_Eu$, dan juga
     
\Centerline{$|h(u)|\le\sup_{E\in\Cal F}|\int_Eu|\le\int|u|$}
     
\noindent untuk setiap $u$, sehingga $h\in(L^1)^*$.
Karena kita mengandaikan bahwa $T$ surjektif, terdapat
$v\in L^{\infty}$ sedemikian sehingga $Tv=h$. Nyatakan $v$ sebagai
$g^{\ssbullet}$ dengan $g:X\to\Bbb R$ terukur dan terbatas secara
esensial. Tetapkan $G=\{x:g(x)>0\}\in\Sigma$.
     
Jika $F\in\Sigma$ dan $\mu F<\infty$, maka
     
\Centerline{$\int_Fg=\int(\chi F)^{\ssbullet}\times g^{\ssbullet}
=(Tv)(\chi F)^{\ssbullet}
=h(\chi F)^{\ssbullet}
=\sup_{E\in\Cal F}\mu(E\cap F)$.}
\Quer\ Jika $E\in\Cal E$ dan $E\setminus G$ tidak terabaikan, terdapat
suatu himpunan $F\subseteq E\setminus G$ sedemikian sehingga
$0<\mu F<\infty$; sekarang
     
     
\Centerline{$\mu F=\mu(E\cap F)\le\int_Fg\le 0$,}
     
\noindent karena $g(x)\le 0$ untuk $x\in F$.\ \BanG\ Jadi
$E\setminus G$ terabaikan untuk setiap $E\in\Cal E$.
     
Misalkan $H\in\Sigma$ sedemikian sehingga $E\setminus H$ terabaikan
untuk setiap $E\in\Cal E$. \Quer\ Jika $G\setminus H$ tidak terabaikan,
terdapat suatu himpunan $F\subseteq G\setminus H$ yang berukuran hingga
dan tak nol. Sekarang
     
     
\Centerline{$\mu(E\cap F)\le\mu(H\cap F)=0$}
     
\noindent untuk setiap $E\in\Cal E$, sehingga
$\mu(E\cap F)=0$ untuk setiap $E\in\Cal F$, dan $\int_Fg=0$; tetapi
$g(x)>0$ untuk setiap $x\in F$, sehingga $\mu F=0$, yang mustahil.
\ \BanG\ Jadi $G\setminus H$ terabaikan.
     
Dengan demikian $G$ merupakan supremum esensial dari $\Cal E$ dalam
$\Sigma$. Karena $\Cal E$ sembarang, $(X,\Sigma,\mu)$ dapat dilokalkan.
     
\medskip
     
\quad{\bf (iii)} Untuk sisa bukti ini, saya akan mengandaikan bahwa
$(X,\Sigma,\mu)$ dapat dilokalkan dan berusaha menunjukkan bahwa $T$
surjektif.
     
Ambil $h\in(L^1)^*$ sedemikian sehingga $\|h\|=1$. Tuliskan
$\Sigma^f=\{F:F\in\Sigma,\,\mu F<\infty\}$, dan untuk
$F\in\Sigma^f$, definisikan $\nu_F:\Sigma\to\Bbb R$ dengan menetapkan
     
\Centerline{$\nu_FE=h(\chi(E\cap F)^{\ssbullet})$}
     
\noindent untuk setiap $E\in\Sigma$. Maka
$\nu_F\emptyset=h(0)=0$, dan jika $E$, $E'\in\Sigma$ saling lepas,

$$\eqalign{\nu_FE+\nu_FE'
&=h(\chi(E\cap F)^{\ssbullet})+h(\chi(E'\cap F)^{\ssbullet})
=h((\chi(E\cap F)+\chi(E'\cap F))^{\ssbullet})\cr
&=h(\chi((E\cup E')\cap F)^{\ssbullet})
=\nu_F(E\cup E').\cr}$$

\noindent Jadi $\nu_F$ aditif. Selain itu,
     
\Centerline{$|\nu_FE|\le\|\chi(E\cap F)^{\ssbullet}\|_1
=\mu(E\cap F)$}

\noindent untuk setiap $E\in\Sigma$, sehingga $\nu_F$ benar-benar
kontinu dalam arti 232Ab. Menurut teorema Radon-Nikod\'ym (232E), terdapat
fungsi terintegralkan $g_F$ sedemikian sehingga $\int_Eg_F=\nu_FE$ untuk
setiap $E\in\Sigma$; kita dapat mengambil $g_F$ terukur dan mempunyai
domain $X$ (232He).
     
\medskip
     
\quad{\bf (iv)} Patut diperhatikan bahwa $|g_F|\le 1$ hampir di mana-mana.
\prooflet{\Prf\ Jika $G=\{x:g_F(x)>1\}$, maka
     
\Centerline{$\int_Gg_F=\nu_FG\le\mu(F\cap G)\le\mu G$;}
     
\noindent tetapi ini hanya mungkin jika $\mu G=0$. Demikian pula, jika
$G'=\{x:g_F(x)<-1\}$, maka
     
     
\Centerline{$\int_{G'}g_F=\nu_FG'\ge-\mu G'$,}
     
\noindent sehingga sekali lagi $\mu G'=0$.\ \Qed}
     
\medskip
     
\quad{\bf (v)} Jika $F$, $F'\in\Sigma^f$, maka $g_F=g_{F'}$ hampir di
mana-mana dalam $F\cap F'$. \Prf\ Jika $E\in\Sigma$ dan
$E\subseteq F\cap F'$, maka
     
\Centerline{$\int_Eg_F=h(\chi(E\cap F)^{\ssbullet})
=h(\chi(E\cap F')^{\ssbullet})=\int_Eg_{F'}$.}
     
\noindent Jadi 131Hb memberikan hasil tersebut.\ \QeD\ 213N (diterapkan
pada $\{g_F\restr F:F\in\Sigma^f\}$) kini memberi tahu kita bahwa, karena
$\mu$ dapat dilokalkan, terdapat fungsi terukur $g:X\to\Bbb R$
sedemikian sehingga $g=g_F$ hampir di mana-mana dalam $F$, untuk setiap
$F\in\Sigma^f$.
     
\medskip
     
\quad{\bf (vi)} Untuk sebarang $F\in\Sigma^f$, himpunan
     
\Centerline{$\{x:x\in F,\,|g(x)|>1\}\subseteq\{x:|g_F(x)|>1\}
\cup\{x:x\in F,\,g(x)\ne g_F(x)\}$}
     
\noindent terabaikan; karena $\mu$ semihingga,
$\{x:|g(x)|>1\}$ terabaikan, dan $g\in\eusm L^{\infty}$, dengan
$\esssup|g|\le 1$. Dengan demikian $v=g^{\ssbullet}\in L^{\infty}$,
dan kita dapat berbicara tentang $Tv\in(L^1)^*$.
     
\medskip
     
\quad{\bf (vii)} Jika $F\in\Sigma^f$, maka

\Centerline{$\int_Fg=\int_Fg_F=\nu_FX=h(\chi F^{\ssbullet})$.}
     
\noindent Segera diperoleh bahwa

\Centerline{$(Tv)(f^{\ssbullet})=\int f\times g=h(f^{\ssbullet})$}
     
\noindent untuk setiap fungsi sederhana $f:X\to\Bbb R$.
Akibatnya $Tv=h$, karena baik $Tv$ maupun $h$ kontinu dan kelas-kelas
ekuivalensi fungsi sederhana membentuk subhimpunan padat dari $L^1$
(242Mb, 2A3Uc). Jadi $h=Tv$ merupakan suatu nilai $T$.
     
\medskip
     
\quad{\bf (viii)} Argumen sebagaimana ditulis di atas mengandaikan bahwa
$\|h\|=1$. Namun tentu saja setiap anggota tak nol dari $(L^1)^*$
merupakan kelipatan skalar dari suatu unsur bernorma $1$, sehingga
merupakan suatu nilai $T$. Jadi $T:L^{\infty}\to(L^1)^*$ memang
surjektif, dan karena itu merupakan isomorfisme isometrik, seperti yang
dinyatakan.
}%end of proof of 243G
     
\leader{243H}{}\cmmnt{ Ingat bahwa $L^0$ selalu lengkap-$\sigma$
Dedekind dan kadang-kadang lengkap Dedekind (241G), sedangkan $L^1$
selalu lengkap Dedekind (242H). Dalam hal ini $L^{\infty}$ mengikuti
$L^0$.
     
\medskip
     
\noindent}{\bf Teorema} Misalkan $(X,\Sigma,\mu)$ suatu ruang ukur.
     
(a) $L^{\infty}(\mu)$ lengkap-$\sigma$ Dedekind.
     
(b) Jika $\mu$ dapat dilokalkan, $L^{\infty}(\mu)$ lengkap Dedekind.
     
\proof{ Keduanya merupakan akibat 241G.
Jika $A\subseteq L^{\infty}=L^{\infty}(\mu)$ terbatas di atas dalam
$L^{\infty}$, tetapkan $u_0\in A$ dan suatu batas atas $w_0$ dari $A$
dalam $L^{\infty}$. Jika $B$ himpunan batas atas untuk $A$ dalam
$L^0=L^0(\mu)$, maka $B\cap L^{\infty}$ merupakan himpunan batas atas
untuk $A$ dalam $L^{\infty}$. Selain itu, jika $B$ mempunyai anggota
terkecil $v_0$, kita harus mempunyai $u_0\le v_0\le w_0$, sehingga
     
\Centerline{$0\le v_0-u_0\le w_0-u_0\in L^{\infty}$}
     
\noindent dan $v_0-u_0$, $v_0$ termasuk dalam $L^{\infty}$.
(Bandingkan bagian (a) dari bukti 242H.) Jadi $v_0=\sup A$ dalam
$L^{\infty}$.
     
Sekarang kita mengetahui bahwa $L^0$ lengkap-$\sigma$ Dedekind; jika
$A\subseteq L^{\infty}$ merupakan himpunan tak kosong terhitung yang
terbatas di atas dalam $L^{\infty}$, tentu himpunan itu terbatas di atas
dalam $L^0$, sehingga mempunyai supremum dalam $L^0$ yang juga merupakan
supremumnya dalam $L^{\infty}$. Karena $A$ sembarang, $L^{\infty}$
lengkap-$\sigma$ Dedekind. Sedangkan jika $\mu$ dapat dilokalkan, kita
dapat berargumen dengan cara yang sama untuk subhimpunan tak kosong
sembarang dari $L^{\infty}$ untuk melihat bahwa $L^{\infty}$ lengkap
Dedekind karena $L^0$ demikian.
}
     
\leader{243I}{Suatu subruang padat dari
\dvrocolon{$L^{\infty}$}}\cmmnt{ Dalam 242M dan 242O saya menguraikan
sepasang subruang linear padat yang penting dari ruang-ruang $L^1$.
Keadaan untuk $L^{\infty}$ agak berbeda. Namun saya dapat menguraikan
satu subruang padat yang penting.
     
\medskip
     
\noindent}{\bf Proposisi} Misalkan $(X,\Sigma,\mu)$ suatu ruang ukur.
     
(a) Tuliskan $\eusm S$ untuk ruang fungsi `sederhana-$\Sigma$' pada $X$,
yaitu ruang fungsi dari $X$ ke $\Bbb R$ yang dapat dinyatakan sebagai
$\sum_{k=0}^na_k\chi E_k$, dengan $a_k\in\Bbb R$ dan $E_k\in\Sigma$
untuk setiap $k\le n$. Maka untuk setiap
$f\in\eusm L^{\infty}=\eusm L^{\infty}(\mu)$ dan setiap $\epsilon>0$,
terdapat $g\in\eusm S$ sedemikian sehingga
$\esssup|f-g|\le\epsilon$.
     
(b) $S=\{f^{\ssbullet}:f\in\eusm S\}$ merupakan subruang linear yang
padat-$\|\,\|_{\infty}$ dalam $L^{\infty}=L^{\infty}(\mu)$.
     
(c) Jika $(X,\Sigma,\mu)$ hingga total, maka $\eusm S$ merupakan ruang
fungsi sederhana-$\mu$, sehingga $S$ hanyalah ruang kelas ekuivalensi
fungsi sederhana\cmmnt{, seperti dalam 242Mb}.
     
\proof{{\bf (a)} Misalkan $\tilde f:X\to\Bbb R$ suatu fungsi terukur
terbatas sedemikian sehingga $f\eae\tilde f$. Misalkan $n\in\Bbb N$
sedemikian sehingga $|f(x)|\le n\epsilon$ untuk setiap $x\in X$.
Untuk $-n\le k\le n$, tetapkan
     
\Centerline{$E_k=\{x:k\epsilon\le \tilde f(x)<(k+1)\epsilon\}$.}
     
\noindent Tetapkan
     
\Centerline{$g=\sum_{k=-n}^nk\epsilon\chi E_k\in\eusm S$;}
     
\noindent maka $0\le\tilde f(x)-g(x)\le\epsilon$ untuk setiap $x\in X$,
sehingga
     
\Centerline{$\esssup|f-g|=\esssup|\tilde f-g|\le\epsilon$.}
     
\medskip
     
{\bf (b)} Ini segera diperoleh, seperti dalam 242Mb.
     
\medskip
     
{\bf (c)} juga elementer.
}%end of proof of 243I
     
     
\vleader{60pt}{243J}{Ekspektasi bersyarat}\cmmnt{ Ekspektasi bersyarat
begitu penting sehingga layak dipertimbangkan interaksinya dengan setiap
konsep baru.

\medskip

}{\bf (a)} Jika $(X,\Sigma,\mu)$ sebarang ruang ukur, dan $\Tau$ suatu
subaljabar-$\sigma$ dari $\Sigma$, maka pembenaman kanonik
$S:L^0(\mu\restrp\Tau)\to L^0(\mu)$\cmmnt{ (242Ja)} membenamkan
$L^{\infty}(\mu\restrp\Tau)$ sebagai subruang dari $L^{\infty}(\mu)$,
dan $\|Su\|_{\infty}=\|u\|_{\infty}$ untuk setiap
$u\in L^{\infty}(\mu\restrp\Tau)$. \dvro{Kita}{Seperti dalam 242Jb,
kita} dapat mengidentifikasi $L^{\infty}(\mu\restrp\Tau)$ dengan
citranya dalam $L^{\infty}(\mu)$.
     
\header{243Jb}{\bf (b)} Sekarang andaikan bahwa $\mu X=1$, dan misalkan
$P:L^1(\mu)\to L^1(\mu\restrp\Tau)$ operator ekspektasi
bersyarat\cmmnt{ (242Jd)}. Maka $L^{\infty}(\mu)$
merupakan\cmmnt{ sebenarnya} subruang linear dari $L^1(\mu)$.
Dengan menetapkan $e=\chi X^{\ssbullet}\in L^{\infty}(\mu)$,\cmmnt{
kita melihat bahwa $\int_Fe=(\mu\restrp\Tau)(F)$ untuk setiap
$F\in\Tau$, sehingga}
     
\Centerline{$Pe=\chi X^{\ssbullet}\in L^{\infty}(\mu\restrp\Tau)$.}
     
\noindent\cmmnt{Jika $u\in L^{\infty}(\mu)$, maka dengan menetapkan
$M=\|u\|_{\infty}$ kita mempunyai $-Me\le u\le Me$, sehingga
$-MPe\le Pu\le MPe$, karena $P$ mempertahankan urutan (242Je); dengan
demikian $\|Pu\|_{\infty}\le M=\|u\|_{\infty}$.
Jadi }$P\restr L^{\infty}(\mu):L^{\infty}(\mu)\to
L^{\infty}(\mu\restrp\Tau)$ merupakan operator bernorma $1$.

Jika $u\in L^{\infty}(\mu\restrp\Tau)$, maka $Pu=u$; jadi
$P[L^{\infty}]$ adalah seluruh $L^{\infty}(\mu\restrp\Tau)$.
     
\leader{243K}{$L^{\infty}$ kompleks}\cmmnt{ Semua gagasan yang
diperlukan untuk menyesuaikan pekerjaan di atas dengan ruang
$L^{\infty}$ kompleks telah muncul dalam 241J dan 242P.} Misalkan
$\eusm L^{\infty}_{\Bbb C}$ adalah
     
\Centerline{$\{f:f\in\eusm L^0_{\Bbb C},\,\esssup|f|<\infty\}
=\{f:\Real(f)\in\eusm L^{\infty},\,\Imag(f)\in\eusm L^{\infty}\}$.}
     
\noindent Maka
     
\Centerline{$L^{\infty}_{\Bbb C}
=\{f^{\ssbullet}:f\in\eusm L^{\infty}_{\Bbb C}\}
=\{u:u\in L^0_{\Bbb C},\,\Real(u)\in L^{\infty},
  \,\Imag(u)\in L^{\infty}\}$.}
     
\noindent Dengan menetapkan
     
\Centerline{$\|u\|_{\infty}
=\||u|\|_{\infty}$\cmmnt{$=\esssup|f|$ setiap kali $f^{\ssbullet}=u$},}
     
\noindent kita mempunyai suatu norma pada $L^{\infty}_{\Bbb C}$ yang
menjadikannya ruang Banach. \cmmnt{Kita tetap mempunyai}
$u\times v\in L^{\infty}_{\Bbb C}$ dan
$\|u\times v\|_{\infty}\le\|u\|_{\infty}\|v\|_{\infty}$ untuk semua
$u$, $v\in L^{\infty}_{\Bbb C}$.
     
Sekarang kita mempunyai dualitas antara $L^1_{\Bbb C}$ dan
$L^{\infty}_{\Bbb C}$ yang menghasilkan suatu operator linear
$T:L^{\infty}_{\Bbb C}\to(L^1_{\Bbb C})^*$ bernorma paling besar $1$,
yang didefinisikan oleh rumus
     
\Centerline{$(Tv)(u)=\int u\times v$
untuk setiap $u\in L^1$, $v\in L^{\infty}$.}
     
\noindent $T$ injektif jika dan hanya jika ruang ukur yang mendasarinya
semihingga, dan bijektif jika dan hanya jika ruang ukur yang mendasarinya
dapat dilokalkan. \prooflet{(Tentu saja ini dapat dibuktikan dengan
mengerjakan kembali argumen 243G; tetapi mungkin lebih mudah untuk
memperhatikan bahwa $T(\Real(v))=\Real(Tv)$,
$T(\Imag(v))=\Imag(Tv)$ untuk setiap $v$, sehingga hasil untuk ruang
kompleks dapat diturunkan dari hasil untuk ruang real.)} \cmmnt{Untuk
memeriksa bahwa} $T$ mempertahankan norma ketika injektif\cmmnt{, rute
tercepat tampaknya meniru argumen (a-ii) dalam bukti 243G}.
     
\exercises{
\leader{243X}{Latihan dasar (a)}
%\spheader 243Xa
Misalkan $(X,\Sigma,\mu)$ sebarang ruang ukur, dan $\hat\mu$
pelengkapan $\mu$ (212C, 241Xb). Tunjukkan bahwa
$\eusm L^{\infty}(\hat\mu)=\eusm L^{\infty}(\mu)$ dan
$L^{\infty}(\hat\mu)=L^{\infty}(\mu)$.
%243A
     
\sqheader 243Xb Misalkan $(X,\Sigma,\mu)$ suatu ruang ukur tak kosong.
Tuliskan $\eusm L^{\infty}_{\Sigma}$ untuk ruang fungsi bernilai real
terbatas yang terukur terhadap $\Sigma$ dan berdomain $X$.
(i) Tunjukkan bahwa $L^{\infty}(\mu)
=\{f^{\ssbullet}:f\in\eusm L^{\infty}_{\Sigma}\}
\subseteq L^0=L^0(\mu)$. (ii) Tunjukkan bahwa
$\eusm L^{\infty}_{\Sigma}$ merupakan kisi Banach lengkap-$\sigma$
Dedekind jika kita memberinya norma
     
\Centerline{$\|f\|_{\infty}=\sup_{x\in X}|f(x)|$ untuk setiap
$f\in\eusm L^{\infty}_{\Sigma}$.}
     
\noindent (iii) Tunjukkan bahwa untuk setiap
$u\in L^{\infty}=L^{\infty}(\mu)$,
$\|u\|_{\infty}=\min\{\|f\|_{\infty}:
f\in\eusm L^{\infty}_{\Sigma},\,f^{\ssbullet}=u\}$.
%243D
     
\sqheader 243Xc Misalkan $(X,\Sigma,\mu)$ sebarang ruang ukur, dan
$A$ suatu subhimpunan dari $L^{\infty}(\mu)$. Tunjukkan bahwa $A$
terbatas untuk norma $\|\,\|_{\infty}$ jika dan hanya jika himpunan itu terbatas
di atas dan di bawah untuk pengurutan $L^{\infty}$.
%243D
     
\spheader 243Xd Misalkan $(X,\Sigma,\mu)$ sebarang ruang ukur, dan
$A\subseteq L^{\infty}(\mu)$ suatu himpunan tak kosong dengan batas atas
terkecil $w$ dalam $L^{\infty}(\mu)$. Tunjukkan bahwa
$\|w\|_{\infty}\le\sup_{u\in A}\|u\|_{\infty}$.
%243D
     
\spheader 243Xe Misalkan
$\langle(X_i,\Sigma_i,\mu_i)\rangle_{i\in I}$ suatu keluarga ruang ukur,
dan $(X,\Sigma,\mu)$ jumlah langsungnya (214L). Tunjukkan bahwa
isomorfisme kanonik antara $L^0(\mu)$ dan
$\prod_{i\in I}L^0(\mu_i)$ (241Xd) menginduksi isomorfisme antara
$L^{\infty}(\mu)$ dan subruang
     
\Centerline{$\{u:u\in\prod_{i\in I}L^{\infty}(\mu_i),\,
\|u\|=\sup_{i\in I}\|u(i)\|_{\infty}<\infty\}$}
     
\noindent dari $\prod_{i\in I}L^{\infty}(\mu_i)$.
%243D
     
\spheader 243Xf Misalkan $(X,\Sigma,\mu)$ sebarang ruang ukur, dan
$u\in L^1(\mu)$. Tunjukkan bahwa terdapat $v\in L^{\infty}(\mu)$
sedemikian sehingga $\|v\|_{\infty}\le 1$ dan
$\int u\times v=\|u\|_1$.
%243F
     
\spheader 243Xg Misalkan $(X,\Sigma,\mu)$ suatu ruang ukur semihingga dan
$v\in L^{\infty}(\mu)$. Tunjukkan bahwa
     
\Centerline{$\|v\|_{\infty}
=\sup\{\int u\times v:u\in L^1,\,\|u\|_1\le 1\}
=\sup\{\|u\times v\|_1:u\in L^1,\,\|u\|_1\le 1\}$.}
%243F

\spheader 243Xh Berikan contoh suatu ruang probabilitas
$(X,\Sigma,\mu)$ dan suatu $v\in L^{\infty}(\mu)$ sedemikian sehingga
$\|u\times v\|_1<\|v\|_{\infty}$ setiap kali $u\in L^1(\mu)$ dan
$\|u\|_1\le 1$.
%243F
     
\spheader 243Xi Tuliskan bukti-bukti 243G yang disesuaikan dengan kasus
khusus (i) $\mu X=1$ (ii) $(X,\Sigma,\mu)$ bersifat $\sigma$-hingga.
%243G
     
\spheader 243Xj Misalkan $(X,\Sigma,\mu)$ sebarang ruang ukur.
Tunjukkan bahwa $L^0(\mu)$ lengkap Dedekind jika dan hanya jika
$L^{\infty}(\mu)$ lengkap Dedekind.
%243G
     
\spheader 243Xk Misalkan $(X,\Sigma,\mu)$ suatu ruang ukur hingga total
dan $\nu:\Sigma\to\Bbb R$ suatu fungsional. Tunjukkan bahwa pernyataan
berikut ekuivalen: (i) terdapat fungsional linear kontinu
$h:L^1(\mu)\to\Bbb R$ sedemikian sehingga
$h((\chi E)^{\ssbullet})=\nu E$ untuk setiap $E\in\Sigma$ (ii) $\nu$
aditif dan terdapat $M\ge 0$ sedemikian sehingga
$|\nu E|\le M\mu E$ untuk setiap $E\in\Sigma$.
%243G
     
\sqheader 243Xl Misalkan $X$ sebarang himpunan, dan misalkan $\mu$ ukuran
pencacahan pada $X$. Dalam kasus ini lazim dituliskan $\ell^{\infty}(X)$
untuk $\eusm L^{\infty}(\mu)$, dan ruang itu diidentifikasi dengan
$L^{\infty}(\mu)$. Tuliskan pernyataan dan bukti hasil-hasil bab ini yang
disesuaikan dengan kasus khusus ini -- jika Anda suka, dengan $X=\Bbb N$.
Secara khusus, tuliskan bukti langsung bahwa $(\ell^1)^*$ dapat
diidentifikasi dengan $\ell^{\infty}$. Apa yang terjadi ketika $X$
hanya mempunyai dua anggota? atau tiga?
%243H
     
\spheader 243Xm Tunjukkan bahwa jika $(X,\Sigma,\mu)$ sebarang ruang ukur
dan $u\in L^{\infty}_{\Bbb C}(\mu)$, maka
     
\Centerline{$\|u\|_{\infty}=\sup\{\|\Real(\zeta u)\|_{\infty}:
\zeta\in\Bbb C,\,|\zeta|=1\}$.}
%243K
     
\spheader 243Xn Misalkan $(X,\Sigma,\mu)$ dan $(Y,\Tau,\nu)$ ruang-ruang
ukur, dan $\phi:X\to Y$ suatu fungsi \imp. Tunjukkan bahwa
$g\phi\in\eusm L^{\infty}(\mu)$ untuk setiap
$g\in\eusm L^{\infty}(\nu)$, dan bahwa peta $g\mapsto g\phi$ menginduksi
suatu operator linear $T:L^{\infty}(\nu)\to L^{\infty}(\mu)$ yang
didefinisikan dengan menetapkan
$T(g^{\ssbullet})=(g\phi)^{\ssbullet}$ untuk setiap
$g\in\eusm L^{\infty}(\nu)$. (Bandingkan 241Xg.) Tunjukkan bahwa
$\|Tv\|_{\infty}=\|v\|_{\infty}$ untuk setiap
$v\in L^{\infty}(\nu)$.
     
\vspheader{48pt}243Xo Untuk $f$, $g\in C=C([0,1])$, ruang fungsi
bernilai real yang kontinu pada interval satuan $[0,1]$, katakanlah
     
\Centerline{$f\le g$ jika dan hanya jika $f(x)\le g(x)$ untuk setiap
$x\in [0,1]$,}
     
\Centerline{$\|f\|_{\infty}=\sup_{x\in [0,1]}|f(x)|$.}
     
\noindent Tunjukkan bahwa $C$ merupakan kisi Banach, dan terlebih lagi
     
\Centerline{$\|f\vee g\|_{\infty}=\max(\|f\|_{\infty},\|g\|_{\infty})$
setiap kali $f$, $g\ge 0$,}
     
\Centerline{$\|f\times g\|_{\infty}\le\|f\|_{\infty}\|g\|_{\infty}$
untuk semua $f$, $g\in C$,}
     
\Centerline{$\|f\|_{\infty}=\min\{\gamma:|f|\le\gamma\chi X\}$ untuk
setiap $f\in C$.}
%243K+
     
\leader{243Y}{Latihan lanjutan (a)} Misalkan $(X,\Sigma,\mu)$ suatu
ruang ukur, dan $Y$ subhimpunan dari $X$; tuliskan $\mu_Y$ untuk ukuran
subruang pada $Y$. Tunjukkan bahwa peta kanonik dari $L^0(\mu)$ pada
$L^0(\mu_Y)$ (241Yg) menginduksi peta kanonik dari
$L^{\infty}(\mu)$ pada $L^{\infty}(\mu_Y)$, yang mempertahankan norma
jika dan hanya jika peta itu injektif jika dan hanya jika $Y$ mempunyai
ukuran luar penuh.
     
}%end of exercises
     
\endnotes{
\Notesheader{243} Saya menyebutkan rumus
     
\Centerline{$\|u\vee v\|_{\infty}=\max(\|u\|_{\infty},\|v\|_{\infty})$
untuk $u$, $v\ge 0$}
     
\noindent (243Df) karena walaupun rumus itu tidak mencirikan ruang
$L^{\infty}$ di antara kisi Banach (lihat 243Xo), dalam suatu arti rumus
tersebut dual terhadap sifat karakteristik
     
\Centerline{$\|u+v\|_1=\|u\|_1+\|v\|_1$ untuk $u$, $v\ge 0$}
     
\noindent dari norma $L^1$. (Saya akan kembali ke hal ini dalam Bab 35
di jilid berikutnya.)
     
Himpunan khusus $\eusm L^{\infty}$ yang saya pilih (243A) agak
sembarang. Ruang $L^{\infty}$ dapat dengan baik diuraikan seluruhnya
sebagai subruang dari $L^0$, tanpa kembali ke fungsi sama sekali; lihat
243Bc, 243Dc. Seperti halnya $\eusm L^0$ dan $\eusm L^1$, ada kesempatan
ketika lebih sederhana untuk bekerja dengan ruang linear fungsi terukur
yang terbatas secara esensial dari $X$ ke $\Bbb R$; dan sekarang kita
mempunyai kandidat jelas yang ketiga, yakni ruang linear
$\eusm L^{\infty}_{\Sigma}$ dari fungsi terukur dari $X$ ke $\Bbb R$ yang
benar-benar terbatas, bukan sekadar terbatas secara esensial, yang
sendirinya merupakan kisi Banach (243Xb).
     
Saya kira teorema terpenting dalam bagian ini adalah 243G, yang
mengidentifikasi $L^{\infty}$ dengan $(L^1)^*$. Identifikasi ini merupakan
alasan utama untuk membedakan ruang ukur `yang dapat dilokalkan'. Bukti
243Gb panjang karena bergantung pada dua gagasan terpisah. Teorema
Radon-Nikod\'ym, pada dasarnya, menangani kasus hingga total, lalu dalam
bagian (b-v) dan (b-vi) bukti tersebut keterlokalan digunakan untuk
merangkai solusi parsial $g_F$. Latihan 243Xi dimaksudkan untuk membantu
Anda membedakan kedua operasi tersebut. Peta
$T':L^1\to(L^{\infty})^*$ (243Fc) juga sangat menarik dengan caranya
sendiri, tetapi akan saya tinggalkan untuk Bab 36.
     
243G memberikan cara lain untuk memandang operator ekspektasi bersyarat.
Jika $(X,\Sigma,\mu)$ suatu ruang probabilitas dan $\Tau$ suatu
subaljabar-$\sigma$ dari $\Sigma$, tentu baik $\mu$ maupun
$\mu\restrp\Tau$ dapat dilokalkan, sehingga $L^{\infty}(\mu)$ dapat
diidentifikasi dengan $(L^1(\mu))^*$ dan
$L^{\infty}(\mu\restrp\Tau)$ dapat diidentifikasi dengan
$(L^1(\mu\restrp\Tau))^*$. Sekarang kita mempunyai pembenaman kanonik
$S:L^1(\mu\restrp\Tau)\to L^1(\mu)$ (242Jb), yang merupakan operator
linear yang mempertahankan norma, sehingga menghasilkan operator adjoint
$S':L^1(\mu)^*\to L^1(\mu\restrp\Tau)^*$ yang didefinisikan oleh rumus
     
\Centerline{$(S'h)(v)=h(Sv)$ untuk semua $v\in L^1(\mu\restrp\Tau)$,
$h\in L^1(\mu)^*$.}
     
\noindent Dengan menuliskan
$T_{\mu}:L^{\infty}(\mu)\to L^1(\mu)^*$ dan
$T_{\mu\restrp\Tau}:L^{\infty}(\mu\restrp\Tau)\to
L^1(\mu\restrp\Tau)^*$ untuk peta-peta kanonik, kita memperoleh peta
$Q=T_{\mu\restrp\Tau}^{-1}S'T_{\mu}:
L^{\infty}(\mu)\to L^{\infty}(\mu\restrp\Tau)$, yang didefinisikan
dengan mengatakan bahwa
     
\Centerline{$\int Qu\times v=(T_{\mu\restrp\Tau}Qu)(v)
=(S'T_{\mu}u)(v)=(T_{\mu}v)(Su)=\int Su\times v=\int u\times v$}
     
\noindent setiap kali $v\in L^1(\mu\restrp\Tau)$ dan
$u\in L^{\infty}(\mu)$. Namun ini sesuai dengan rumus 242L: kita
mempunyai
     
\Centerline{$\int Qu\times v=\int u\times v=\int P(u\times v)
=\int Pu\times v$.}
     
\noindent Karena $v$ sembarang, kita harus mempunyai $Qu=Pu$ untuk
setiap $u\in L^{\infty}(\mu)$. Jadi, dalam suatu arti, operator
ekspektasi bersyarat merupakan adjoint dari operator pembenaman yang
sesuai.
     
Pembahasan dalam paragraf terakhir tentu saja hanya berlaku bagi
pembatasan $P\restr L^{\infty}(\mu)$ dari operator ekspektasi bersyarat
pada ruang $L^{\infty}$. Karena $\mu$ hingga total,
$L^{\infty}(\mu)$ merupakan subruang dari $L^1(\mu)$, dan sifat-sifat
hakiki operator $P$ berhubungan dengan perilakunya pada seluruh ruang
$L^1$. $P:L^1(\mu)\to L^1(\mu\restrp\Tau)$ juga dapat dinyatakan sebagai
operator adjoint, tetapi ungkapannya memerlukan lebih banyak teori ruang
Riesz daripada yang dapat saya muat di sini. Saya akan kembali ke topik
ini dalam Bab 36.
}%end of notes
     
\discrpage
 
